% !TeX root = main.tex
% PAGINE ORIGINALI 421-430 — DA COMPILARE
% Incollare qui i quiz del blocco e alla fine le rispettive\registraSoluzione{numero}{lettera}.
% Questo file vuoto non aggiunge pagine al PDF.

% !TeX root = main.tex
% Pagine originali 421-430: 19a simulazione (3 quesiti), 20a simulazione (20), 21a simulazione (5).

% Pag. originale 421/521
\iffalse Pag. originale 421/521 \fi
\begin{quiz}
$\displaystyle\lim_{x\to0^-}\frac{\sin(-1/x)-4}{x}=$
\begin{enumerate}
\item $0$
\item $+\infty$
\item $-\infty$
\item non esiste
\item $-4$
\end{enumerate}
\end{quiz}
% SOLUZIONE: B
% SEZIONE_AM1: 3.1
% SPIEGAZIONE: Il numeratore \`e compreso fra $-5$ e $-3$, mentre $x\to0^-$: il rapporto tende a $+\infty$.

% Pag. originale 421/521
\iffalse Pag. originale 421/521 \fi
\begin{quiz}
Sia $z=-2-i$, qual \`e il modulo del numero complesso $z^2+3i$?
\begin{enumerate}
\item $14$
\item $8$
\item $10$
\item $\sqrt{58}$
\item $\sqrt{74}$
\end{enumerate}
\end{quiz}
% SOLUZIONE: D
% SEZIONE_AM1: 1.5
% SPIEGAZIONE: $z^2=3+4i$, dunque $z^2+3i=3+7i$ e il modulo vale $\sqrt{9+49}=\sqrt{58}$.

% Pag. originale 421/521
\iffalse Pag. originale 421/521 \fi
\begin{quiz}
Una primitiva della funzione $f(x)=2\sin^2x\cdot\cos x$ \`e:
\begin{enumerate}
\item $\frac13\sin^3x+\frac13$
\item $\sin^3x$
\item $-\frac23\sin^3x\cdot\sin x$
\item $-\frac23\sin^3x+2$
\item $\frac23\sin^3x+\sin^2x+\cos^2x$
\end{enumerate}
\end{quiz}
% SOLUZIONE: E
% SEZIONE_AM1: 5.1
% SPIEGAZIONE: Ponendo $t=\sin x$ si ottiene $\int2t^2\,dt=\frac23t^3+C$; poich\'e $\sin^2x+\cos^2x=1$, E \`e una primitiva.

% Pag. originale 424/521
\iffalse Pag. originale 424/521 \fi
\begin{quiz}
Siano date le due funzioni $f(x)=\dfrac{\sqrt{x-1}}{\sqrt{x+3}}$, $g(x)=\sqrt{\dfrac{x-1}{x+3}}$. Quale delle seguenti propriet\`a \`e soddisfatta?
\begin{enumerate}
\item $f(x)=g(x)$, $\forall x\in(-\infty,-3)\cup[1,+\infty)$
\item $f(x)=g(x)$, $\forall x\in(1,+\infty)$
\item $\operatorname{dom}g\subseteq\operatorname{dom}f$
\item $\operatorname{dom}g=\operatorname{dom}f$
\item $f(x)=g(x)$, $\forall x\in[1,+\infty)$
\end{enumerate}
\end{quiz}
% SOLUZIONE: E
% SEZIONE_AM1: 1.4
% SPIEGAZIONE: $\operatorname{dom}f=[1,+\infty)$ e $\operatorname{dom}g=(-\infty,-3)\cup[1,+\infty)$. Le due espressioni coincidono su $[1,+\infty)$. Anche B \`e vera come affermazione non massimale: ambiguit\`a nella fonte.

% Pag. originale 424/521
\iffalse Pag. originale 424/521 \fi
\begin{quiz}
$\displaystyle\int_0^\pi\cos x\,e^{\sin x}\,dx=$
\begin{enumerate}
\item $0$
\item $1$
\item $2$
\item $e^\pi-1$
\item $-1$
\end{enumerate}
\end{quiz}
% SOLUZIONE: A
% SEZIONE_AM1: 5.1
% SPIEGAZIONE: Una primitiva \`e $e^{\sin x}$: agli estremi vale $1-1=0$.

% Pag. originale 424/521
\iffalse Pag. originale 424/521 \fi
\begin{quiz}
Il seguente integrale improprio $\displaystyle\int_0^{+\infty}\frac{\sqrt{x+2}}{x+5}\,dx$
\begin{enumerate}
\item converge a zero
\item diverge a $+\infty$
\item \`e indeterminato
\item converge ad un numero positivo
\item nessuna delle risposte precedenti
\end{enumerate}
\end{quiz}
% SOLUZIONE: B
% SEZIONE_AM1: 5.5
% SPIEGAZIONE: Per $x\to+\infty$ l\'integranda \`e equivalente a $1/\sqrt x$, il cui integrale diverge.

% Pag. originale 424/521
\iffalse Pag. originale 424/521 \fi
\begin{quiz}
Sia data una funzione continua nel suo dominio, tale che $f:\mathbb R\setminus\{1,6\}\to\mathbb R$. Allora ha sicuramente un punto di minimo se:
\begin{enumerate}
\item $\lim_{x\to1^-}f(x)=+\infty$ e $\lim_{x\to6^+}f(x)=+\infty$
\item $\lim_{x\to1^+}f(x)=-\infty$ e $\lim_{x\to6^-}f(x)=+\infty$
\item $\inf(f)\leq0$; $\sup(f)\geq0$
\item $\inf(f)\leq0$; $\sup(f)>0$
\item $\lim_{x\to1^+}f(x)=+\infty$ e $\lim_{x\to6^-}f(x)=+\infty$
\end{enumerate}
\end{quiz}
% SOLUZIONE: E
% SEZIONE_AM1: 4.3
% SPIEGAZIONE: Sul tratto $(1,6)$ la funzione continua tende a $+\infty$ presso entrambi gli estremi; quindi raggiunge un minimo interno.

% Pag. originale 424/521
\iffalse Pag. originale 424/521 \fi
\begin{quiz}
Sia $A\subseteq\mathbb R$ e sia $p=\inf A$. Sapendo che $p\in\mathbb R$, quale delle seguenti affermazioni \`e FALSA?
\begin{enumerate}
\item $A$ \`e limitato inferiormente
\item se $p\in A\Rightarrow p$ \`e il minimo di $A$
\item $p$ \`e il pi\`u grande dei minoranti
\item $A$ \`e limitato
\item $p$ \`e un minorante di $A$
\end{enumerate}
\end{quiz}
% SOLUZIONE: D
% SEZIONE_AM1: 1.5
% SPIEGAZIONE: Un estremo inferiore finito garantisce soltanto limitatezza inferiore, non superiore: ad esempio $A=[0,+\infty)$.

% Pag. originale 425/521
\iffalse Pag. originale 425/521 \fi
\begin{quiz}
Sia $f$ una funzione continua in $[10,+\infty)$ che soddisfa alla condizione $f(x)\geq\frac1{x^{10}}$, $\forall x\geq10$. Allora $\displaystyle\int_{10}^{+\infty}f(x)\,dx$
\begin{enumerate}
\item converge
\item \`e indeterminato
\item diverge
\item pu\`o convergere (ad un numero positivo) oppure divergere a $+\infty$
\item converge a zero
\end{enumerate}
\end{quiz}
% SOLUZIONE: D
% SEZIONE_AM1: 5.5
% SPIEGAZIONE: Il minorante $x^{-10}$ \`e integrabile, ma ci\`o non decide la convergenza di $f$: scegliere $f=x^{-10}$ oppure $f=1$.

% Pag. originale 425/521
\iffalse Pag. originale 425/521 \fi
\begin{quiz}
Quale delle seguenti funzioni NON \`e la derivata della funzione $f(x)=\log\bigl(\cos(2x)e^{3x}\bigr)$?
\begin{enumerate}
\item $f\prime(x)=\dfrac{3\cos2x-2\sin2x}{\cos2x}$
\item $f\prime(x)=\dfrac{3\cos2x-4\sin x\cos x}{\cos2x}$
\item $f\prime(x)=3-2\tan2x$
\item $f\prime(x)=3-2\dfrac{\sin2x}{\cos2x}$
\item $f\prime(x)=3-\dfrac{\sin2x}{\cos2x}$
\end{enumerate}
\end{quiz}
% SOLUZIONE: E
% SEZIONE_AM1: 4.1
% SPIEGAZIONE: La derivata \`e $3-2\tan(2x)$; l\'opzione E manca del fattore $2$.

% Pag. originale 425/521
\iffalse Pag. originale 425/521 \fi
\begin{quiz}
L\'equazione differenziale $y\prime=\dfrac{x-y}{x+1}$, $x\in(-1,+\infty)$
\begin{enumerate}
\item ha la soluzione $y(x)=x$
\item ha la soluzione $y(x)=1$
\item \`e a variabili separabili
\item \`e lineare
\item non ha soluzioni
\end{enumerate}
\end{quiz}
% SOLUZIONE: D
% SEZIONE_AM1: 6.1
% SPIEGAZIONE: Riscrivendo $y\prime+\frac{1}{x+1}y=\frac{x}{x+1}$ si riconosce un\'equazione lineare del primo ordine.

% Pag. originale 425/521
\iffalse Pag. originale 425/521 \fi
\begin{quiz}
Calcolare il limite $\displaystyle\lim_{x\to+\infty}\frac{e^{\frac{x^2+3}{x-4}}}{e^x}=$
\begin{enumerate}
\item $1$
\item $0$
\item $+\infty$
\item $e$
\item $e^4$
\end{enumerate}
\end{quiz}
% SOLUZIONE: E
% SEZIONE_AM1: 3.1
% SPIEGAZIONE: $\frac{x^2+3}{x-4}-x=\frac{4x+3}{x-4}\to4$; il limite vale $e^4$.

% Pag. originale 425/521
\iffalse Pag. originale 425/521 \fi
\begin{quiz}
$\displaystyle\lim_{x\to0}\frac{1-\cos4x}{\sinh^2(3x)}=$
\begin{enumerate}
\item $\frac{16}9$
\item $\frac89$
\item $\frac23$
\item $\frac43$
\item $\frac29$
\end{enumerate}
\end{quiz}
% SOLUZIONE: B
% SEZIONE_AM1: 3.1
% SPIEGAZIONE: $1-\cos4x\sim8x^2$ e $\sinh^2(3x)\sim9x^2$, quindi $8/9$.

% Pag. originale 426/521
\iffalse Pag. originale 426/521 \fi
\begin{quiz}
La parte principale, per $x\to0$, della funzione $f(x)=\sqrt[7]{x^6+x^7+3x}$, rispetto al campione standard $u(x)=x$, \`e
\begin{enumerate}
\item $3x$
\item $3x^{1/7}$
\item $x$
\item $x^{6/7}$
\item $(3x)^{1/7}$
\end{enumerate}
\end{quiz}
% SOLUZIONE: E
% SEZIONE_AM1: 3.4
% SPIEGAZIONE: All\'interno della radice il termine dominante \`e $3x$; la radice settima \`e definita anche per $x<0$.

% Pag. originale 426/521
\iffalse Pag. originale 426/521 \fi
\begin{quiz}
$\displaystyle\lim_{n\to+\infty}\left(\cos\left(\frac{n}{n^2+1}\right)\right)^{-\arctan n}$
\begin{enumerate}
\item $+\infty$
\item $0$
\item $1$
\item $\frac\pi2$
\item non esiste
\end{enumerate}
\end{quiz}
% SOLUZIONE: C
% SEZIONE_AM1: 2.2
% SPIEGAZIONE: La base tende a $1$ e l\'esponente tende a $-\pi/2$, quindi il limite \`e $1$.

% Pag. originale 426/521
\iffalse Pag. originale 426/521 \fi
\begin{quiz}
Il polinomio di Mac Laurin, di ordine $2$, della funzione $f(x)=\dfrac{2e^{-x}-x^2}{1+x^2}$ \`e
\begin{enumerate}
\item $1-x-x^2$
\item $2x-2x^2$
\item $2-2x-2x^2+o(x^2)$
\item $2-2x-2x^2$
\item $x-x^2$
\end{enumerate}
\end{quiz}
% SOLUZIONE: D
% SEZIONE_AM1: 4.2
% SPIEGAZIONE: $2e^{-x}-x^2=2-2x+o(x^2)$ e $(1+x^2)^{-1}=1-x^2+o(x^2)$, quindi il polinomio \`e $2-2x-2x^2$. Anche C \`e lo sviluppo completo con resto.

% Pag. originale 426/521
\iffalse Pag. originale 426/521 \fi
\begin{quiz}
Sia data la funzione $f:\mathbb R\to\mathbb R$, continua e derivabile e tale che $f(1)=3$, $f(5)=7$. Allora $\exists c\in(1,5)$ tale che:
\begin{enumerate}
\item $f\prime(c)=2$
\item $f\prime(c)=0$
\item $f\prime(c)=1$
\item $f\prime(c)=-2$
\item $f\prime(c)=5$
\end{enumerate}
\end{quiz}
% SOLUZIONE: C
% SEZIONE_AM1: 4.3
% SPIEGAZIONE: Il teorema di Lagrange d\`a $f\prime(c)=(7-3)/(5-1)=1$.

% Pag. originale 426/521
\iffalse Pag. originale 426/521 \fi
\begin{quiz}
Lo sviluppo di Taylor della funzione $f$, nell\'intorno di $x=2$, \`e $f(x)=e+e(x-2)-2e(x-2)^2+o((x-2)^2)$. Allora
\begin{enumerate}
\item $f$ ha minimo in $x=2$
\item $f$ ha massimo in $x=2$
\item $f$ ha un flesso in $x=2$
\item la retta tangente in $x=0$ ha equazione $y=-e+ex$
\item $\exists\delta>0:\forall x\in I_\delta(2)$ la funzione sia positiva, crescente e concava
\end{enumerate}
\end{quiz}
% SOLUZIONE: E
% SEZIONE_AM1: 4.3
% SPIEGAZIONE: Dallo sviluppo: $f(2)=e$, $f\prime(2)=e$, $f\prime\prime(2)=-4e$. L\'affermazione di concavit\`a in un intorno richiede regolarit\`a aggiuntiva (per esempio continuit\`a di $f\prime\prime$).

% Pag. originale 426/521
\iffalse Pag. originale 426/521 \fi
\begin{quiz}
Sia $I\subset\mathbb R$ un intervallo e siano $f,g$ due funzioni definite da $I\to\mathbb R$ che soddisfano alla condizione $\forall x\in I\ f\prime(x)=g\prime(x)$, allora
\begin{enumerate}
\item $f-g$ \`e una funzione costante in $I$
\item $f,g$ hanno la stessa primitiva in $I$
\item $f,g$ hanno integrali indefiniti uguali in $I$
\item $f(x)=g(x)$, $\forall x\in I$
\item $\int_a^b f(x)\,dx=\int_a^b g(x)\,dx$, $\forall a,b\in I$
\end{enumerate}
\end{quiz}
% SOLUZIONE: A
% SEZIONE_AM1: 5.1
% SPIEGAZIONE: Per il teorema di Lagrange $(f-g)\prime=0$ su un intervallo implica $f-g=C$.

% Pag. originale 427/521
\iffalse Pag. originale 427/521 \fi
\begin{quiz}
Siano $f,g$ due funzioni definite da $\mathbb R\to\mathbb R$, entrambe derivabili, con $f\prime(0)=1$, $f(0)=-3$. Posto $h(x)=f(x)\cdot g(x)$ allora:
\begin{enumerate}
\item $h\prime(0)=-3g\prime(0)+g(0)$
\item $h\prime(0)=g\prime(0)$
\item $h\prime(0)=3g\prime(0)$
\item $h\prime(0)=-3$
\item $h\prime(3)=-3g(3)+g\prime(3)$
\end{enumerate}
\end{quiz}
% SOLUZIONE: A
% SEZIONE_AM1: 4.1
% SPIEGAZIONE: La regola del prodotto fornisce $h\prime(0)=f\prime(0)g(0)+f(0)g\prime(0)=g(0)-3g\prime(0)$.

% Pag. originale 427/521
\iffalse Pag. originale 427/521 \fi
\begin{quiz}
Sia $f(x)$ una funzione che soddisfa alla condizione $x+1<f(x)<x+2$, $\forall x\in\mathbb R$. Allora
\begin{enumerate}
\item $f(x)\sim x$, $x\to+\infty$
\item la retta $y=x+\frac32$ \`e asintoto obliquo per la funzione
\item la funzione \`e monotona, strettamente crescente in $\mathbb R$
\item l\'equazione della funzione $f(x)$ \`e del tipo $f(x)=x+k$, $k\in[1,2]$
\item $f(x)\geq0$, $\forall x\in\mathbb R$
\end{enumerate}
\end{quiz}
% SOLUZIONE: A
% SEZIONE_AM1: 3.4
% SPIEGAZIONE: Dividendo per $x>0$, $1+1/x<f(x)/x<1+2/x$; per confronto $f(x)/x\to1$.

% Pag. originale 427/521
\iffalse Pag. originale 427/521 \fi
\begin{quiz}
La funzione $f(x)=|x-4|^a$, dove $a$ \`e un parametro positivo, \`e derivabile due volte in $\mathbb R$ se e solo se
\begin{enumerate}
\item $a>2$
\item $a>0$
\item $a\geq2$
\item $a\geq1$
\item $a>1$
\end{enumerate}
\end{quiz}
% SOLUZIONE: C
% SEZIONE_AM1: 4.1
% SPIEGAZIONE: Per $a=2$ \`e un polinomio; per $a>2$ la derivata seconda tende a zero nel punto $4$. Per $1<a<2$ la seconda derivata non esiste in $4$.

% Pag. originale 427/521
\iffalse Pag. originale 427/521 \fi
\begin{quiz}
Se $f:\mathbb R\to\mathbb R$ \`e una funzione continua e derivabile almeno $5$ volte nell\'intorno di $x=0$, con sviluppo di Mac Laurin $f(x)=3x^3+x^5+o(x^5)$, allora l\'integrale improprio $\displaystyle\int_0^1\frac{f(x)\sin x}{x^\alpha}\,dx$
\begin{enumerate}
\item converge per ogni $\alpha<1$
\item converge per ogni $\alpha<5$
\item converge per ogni $\alpha>1$
\item converge per ogni $\alpha>5$
\item \`e indeterminato per ogni $\alpha$
\end{enumerate}
\end{quiz}
% SOLUZIONE: B
% SEZIONE_AM1: 5.5
% SPIEGAZIONE: Vicino a zero $f(x)\sin x\sim3x^4$, quindi l\'integranda \`e equivalente a $3x^{4-\alpha}$ e converge se $4-\alpha>-1$. Anche A \`e vera, ma meno precisa: ambiguit\`a nel quesito.

% Pag. originale 430/521
\iffalse Pag. originale 430/521 \fi
\begin{quiz}
Quale delle seguenti classi di funzioni NON \`e sempre integrabile secondo Riemann nell\'intervallo $[a,b]$?
\begin{enumerate}
\item le funzioni continue a tratti su $[a,b]$
\item le funzioni continue su $(a,b)$ e limitate su $[a,b]$
\item le funzioni monotone su $[a,b]$
\item le funzioni limitate su $[a,b]$
\item le funzioni continue su $[a,b]$
\end{enumerate}
\end{quiz}
% SOLUZIONE: D
% SEZIONE_AM1: 5.1
% SPIEGAZIONE: La sola limitatezza non basta: la funzione di Dirichlet, $1$ sui razionali e $0$ sugli irrazionali, non \`e Riemann-integrabile.

% Pag. originale 430/521
\iffalse Pag. originale 430/521 \fi
\begin{quiz}
Quale delle seguenti propriet\`a NON \`e soddisfatta da una funzione Riemann integrabile in $[a,b]$?
\begin{enumerate}
\item condizione sufficiente affinch\'e una funzione $f$, definita in $[a,b]$, sia Riemann integrabile su $[a,b]$ \`e che sia continua in $[a,b]$
\item condizione sufficiente affinch\'e una funzione $f$, definita in $[a,b]$, sia Riemann integrabile su $[a,b]$ \`e che sia derivabile in $[a,b]$
\item condizione necessaria affinch\'e una funzione $f$, definita in $[a,b]$, sia Riemann integrabile su $[a,b]$ \`e che sia continua in $[a,b]$
\item condizione sufficiente affinch\'e una funzione $f$, definita in $[a,b]$, sia Riemann integrabile su $[a,b]$ \`e che sia continua in $(a,b)$ e limitata in $[a,b]$
\item condizione sufficiente affinch\'e una funzione $f$, definita in $[a,b]$, sia Riemann integrabile su $[a,b]$ \`e che sia continua a tratti in $[a,b]$
\end{enumerate}
\end{quiz}
% SOLUZIONE: C
% SEZIONE_AM1: 5.1
% SPIEGAZIONE: La continuit\`a \`e sufficiente ma non necessaria: le funzioni con un numero finito di discontinuit\`a limitate sono integrabili.

% Pag. originale 430/521
\iffalse Pag. originale 430/521 \fi
\begin{quiz}
Quante sono le soluzioni $z\in\mathbb C$ dell\'equazione $(z^4+4)(z^3-8)=0$ con $\operatorname{Re}(z)<0$?
\begin{enumerate}
\item quattro
\item tre
\item due
\item sette
\item nessuna
\end{enumerate}
\end{quiz}
% SOLUZIONE: A
% SEZIONE_AM1: 1.5
% SPIEGAZIONE: Le quattro radici di $z^4=-4$ sono $\pm1\pm i$, di cui due con parte reale negativa. Le tre radici di $z^3=8$ sono $2$ e $-1\pm i\sqrt3$, di cui altre due negative: totale quattro.

% Pag. originale 430/521
\iffalse Pag. originale 430/521 \fi
\begin{quiz}
$\displaystyle\lim_{x\to0^+}\frac{\log(1-\cos(7x^2))}{\log x}=$
\begin{enumerate}
\item $2$
\item $4$
\item $-2$
\item $+\infty$
\item $49$
\end{enumerate}
\end{quiz}
% SOLUZIONE: B
% SEZIONE_AM1: 3.1
% SPIEGAZIONE: $1-\cos(7x^2)\sim(49/2)x^4$, dunque il logaritmo \`e $4\log x+O(1)$ e il limite vale $4$.

% Pag. originale 430/521
\iffalse Pag. originale 430/521 \fi
\begin{quiz}
$\displaystyle\int_{\log2}^{\log10}\frac{e^x}{\sqrt{e^x-1}}\,dx=$
\begin{enumerate}
\item $6-\sqrt3$
\item $2$
\item $4$
\item $2\log10-2\log2$
\item $6$
\end{enumerate}
\end{quiz}
% SOLUZIONE: C
% SEZIONE_AM1: 5.1
% SPIEGAZIONE: La primitiva \`e $2\sqrt{e^x-1}$; il valore \`e $2(3-1)=4$.
Risposta completata
